% AUTORIDAD DE FUENTE: Exposé IX, §6, líneas 1783--1842;
% índices de escaneo 403--410, folios 392--399. Las tres diferencias
% materiales D/C se conservan reversiblemente mediante \sisi. Se han
% normalizado sin aparato varios errores puramente transcriptivos que la
% impresión resuelve: p^\infty, p^N A_o^o, A'_{K'} en (6.3.2), la
% normalización en K' y la función de traza sobre \pi'.

\medskip
\noindent
\underline{6. Observaciones sobre la construcción de la parte tórica y de
la parte fija de} $T_p(A)$ \underline{en el caso de reducción no
semiestable}.

\medskip
\noindent
6.1. La construcción de la parte tórica de $T_\ell(A_K)$ dada en 5.1
sigue teniendo sentido en cuanto $S$ es completo, sin suponer que
$A^o$ sea $\ell$-divisible; es decir (cuando $p=\ell$), sin suponer la
reducción semiestable. Esta hipótesis solo se había utilizado para
formular el teorema de ortogonalidad 5.2. Además, se puede construir
$T_p(A_K)^t$ suponiendo únicamente que $S$ es henseliano: se demuestra
que la parte tórica $T_p(A_{\widehat K})^t$ sobre el completado,
construida por el método de 5.1, proviene de un subprogrupo de
Barsotti--Tate de $T_p(A_K)$, necesariamente determinado de manera
única.

Más precisamente, se puede construir un grupo de Barsotti--Tate
pro-$p$, $T_p(A)^t$, \underline{sobre} $S$, y un homomorfismo
\[
 T_p(A)^t\longrightarrow T_p(A)
\]
de sistemas proyectivos que es un monomorfismo término a término y cuya
imagen inversa sobre $\widehat S$ da los objetos análogos construidos
por el método de 5.1. Esto se deduce parafraseando la construcción de
5.1, salvo que, en vez de levantar todo el toro maximal $T_o$ de $A_o$,
para cada entero $\nu\geq0$ basta levantar el subgrupo
${}_{p^\nu}T_o$ a un subgrupo de tipo multiplicativo
$({}_{p^\nu}A^o)^t$ de ${}_{p^\nu}A^o$.

%%%%%%%%%%  scan idx 404  |  folio 393  %%%%%%%%%%

Para la existencia y unicidad de tal levantamiento, obsérvese que el
funtor
\[
 (\mathrm{Sch})^o_{/S}\longrightarrow\mathrm{Ens},
 \qquad
 S'\longmapsto
 \{\text{subgrupos finitos de tipo multiplicativo de }A^o_{S'}\}
\]
está representado por un $S$-esquema $X$ de tipo finito (cf. la
demostración de SGA 3 XI 3.12 a)), que es étale sobre $S$ por SGA 3 IX
3.6 bis. En consecuencia, existe exactamente una sección de $X$ sobre
$S$ que prolonga una sección dada de $X_s$ sobre $s$, lo cual demuestra
la afirmación.

El mismo argumento demuestra, más generalmente, que, para todo
\underline{subgrupo finito de tipo multiplicativo} $H_o$ de $A_o$,
existe un único subgrupo finito de tipo multiplicativo $H$ de $A$ que
lo prolonga. Por consiguiente, para todo subgrupo $p$ de Barsotti--Tate
\underline{tórico} $H_o$ de $A_o$ (es decir, uno cuyas componentes son
de tipo multiplicativo), existe un único subgrupo $p$ de Barsotti--Tate
$H$ de $A$ que lo prolonga. Esto da a su vez un subgrupo $p$ de
Barsotti--Tate de ${}_{p^\infty}A_K$, o, equivalentemente, un subgrupo
pro-$p$ de Barsotti--Tate de $T_p(A_K)$.

En particular, este procedimiento se aplica cuando $H_o$ es el mayor
subgrupo $p$ de Barsotti--Tate tórico de $A_o$. Se obtiene tomando el
mayor subgrupo $p$-divisible $\overline A_o$ de $A_o^o$
($\overline A_o=p^N A_o^o$ para $N$ suficientemente grande) y la
``parte tórica'' de ${}_{p^\infty}\overline A_o$ (definida, por ejemplo,
por SGA 3 XVII 7.2.1). Obtenemos así un subgrupo de Barsotti--Tate de
${}_{p^\infty}A_K$ que podría llamarse la \underline{parte tórica} de
este último.

\medskip
\noindent
6.2. Sea de nuevo $\overline A_o$ el mayor subesquema en grupos
$p$-divisible de $A_o^o$, y consideremos el grupo $p$ de
Barsotti--Tate
\[
 F_o={}_{p^\infty}(\overline A_o).
\]
Es inmediato que $F_o$ puede caracterizarse como el mayor subgrupo $p$
de Barsotti--Tate de $A_o$ (o de $A_o^o$). Es natural preguntarse si
este grupo de Barsotti--Tate puede prolongarse a un subgrupo de
Barsotti--Tate $F$ de $A$, que desempeñaría entonces el papel de la
\underline{parte fija} $T_\ell(A)^f$ de (2.2.3.2) en el caso en que
$A^o$ es $\ell$-divisible.

%%%%%%%%%%  scan idx 405  |  folio 394  %%%%%%%%%%

En el caso semiestable, donde $\overline A_o=A_o^o$, el método de 2.2
responde afirmativamente a la pregunta. En el caso general, siempre que
el problema tenga solución, se ve fácilmente que esta es única. Los
subprogrupos correspondientes de $T_p(A_K)$ y $T_p(A'_K)$ (que también
merecen llamarse las \underline{partes fijas} de los progrupos en
cuestión) siguen satisfaciendo el teorema de ortogonalidad 5.2, y la
parte fija $T_p(A_K)^f$ todavía puede caracterizarse como en 5.8.

\medskip
\noindent
6.2.1. Demostraremos, sin embargo, que en general la pregunta no admite
una respuesta afirmativa; es decir, \underline{la parte fija
$T_p(A)^f$ no puede definirse en general} en el caso de reducción no
semiestable. En efecto, por un teorema de levantamiento de Serre--Tate
[30], cuando $S$ es completo, la existencia de la parte fija equivale a
la existencia de levantamientos infinitesimales $\overline A_n$ de
$\overline A_o$ en los $A_n$, o también a la existencia de un esquema
en grupos liso $\overline A$ sobre $S$, extensión de un esquema
abeliano por un toro $T$, con fibra especial $\overline A_o$, junto con
un homomorfismo formal
\[
 \widehat{\overline A}\longrightarrow\widehat A
\]
que prolonga la inclusión en las fibras especiales.

Así, cuando $\overline A_o$ es un esquema abeliano, es decir, cuando
$T_o=0$, esto equivale a la existencia de un esquema abeliano
$\overline A$ que prolongue $\overline A_o$ y de un homomorfismo
$\overline A\longrightarrow A$ que prolongue la inclusión de las fibras
especiales. Bastará, pues, dar un ejemplo en el que $T_o=0$, $A_o^o$ no
sea ni unipotente ni propio y $A_K$ sea simple sobre $K$.

%%%%%%%%%%  scan idx 406  |  folio 395  %%%%%%%%%%

\medskip
\noindent
6.2.2. He aquí un ejemplo debido a M. Raynaud. Se basa en la observación
de que, si $S$ es un rasgo, $S'$ es la normalización de $S$ en una
extensión separable finita $K'$ de $K$, y $A'$ es un esquema abeliano
sobre $S'$, entonces el esquema en grupos
\[
 A=\prod_{S'/S} A'/S'
\]
sobre $S$ es un modelo de Néron: se identifica con el modelo de Néron de
su fibra genérica, que es, en efecto, un esquema abeliano. Esto se
deduce inmediatamente de las definiciones y del hecho de que $A'$ es él
mismo un modelo de Néron sobre $S'$. Entonces
\[
 A_o=\prod_{S'_s/s}(A'\times_{S'}S'_s)/S'_s
\]
es una extensión de un esquema abeliano, no nulo si $A'\neq0$, por un
esquema unipotente conexo, no nulo si $A'\neq0$ y $S'$ está ramificado
sobre $S$.

Basta, por tanto, dar un ejemplo en el que $A'\neq0$, $S'$ esté
ramificado sobre $S$ (un rasgo local completo), y
\[
 A_K=\prod_{K'/K}A'_{K'}/K'
\]
sea \sisi{simple sobre $K'$}{simple sobre $K$}. Tomemos $K'$ como una
extensión cuadrática ramificada de \sisi{$S$}{$K$} (tal extensión existe
siempre), y
tomemos $A'$ como una curva elíptica sobre $S'$. Entonces $A$ tiene
dimensión relativa dos, y decir que $A_K$ es simple sobre $K$ significa
que no existe ningún homomorfismo no nulo
\[
 B_K\longrightarrow A_K,
\]
donde $B_K$ es un esquema abeliano de dimensión uno sobre $K$. Dar tal
homomorfismo equivale a dar un homomorfismo no nulo
\[
 B_{K'}\longrightarrow A'_{K'},
\]
que sería necesariamente una isogenia. Por tanto, $A_K$ es simple
sobre $K$ precisamente cuando $A'_{K'}$ no es isógena a una curva
elíptica proveniente de $K$.

Si ocurriera lo contrario, el invariante $j$ de $A'$ dependería
algebraicamente (sobre el cuerpo primo) de su conjugado $\overline j$
bajo el automorfismo no trivial de $K'$ sobre $K$. Si se elige $j$ de
modo que
\[
 a=j+\overline j=\operatorname{Tr}_{K'/K}j
\]
sea trascendente sobre el cuerpo primo, esto equivale a decir que $j$
depende algebraicamente de $a$. Sin embargo, por consideraciones de
cardinalidad, existe un elemento $a\in V$ trascendente sobre el
subcuerpo primo $P$ de $K$, y luego un elemento $b\in V'$ (la
normalización de $V$ en $K'$) trascendente sobre $P(a)$ y tal que
\[
 \operatorname{Tr}_{V'/V}b=0.
\]

%%%%%%%%%%  scan idx 407  |  folio 396  %%%%%%%%%%

Basta tomar $j=a+b\in V'$. Podemos suponer además que
$b\in\mathfrak r(V')$, y que la imagen de $a$ en $k=V/\mathfrak m$ es
un elemento prescrito de $k$, elegido distinto de los valores
excepcionales $j_o=0$ y $j_o=12^3$. Esto es posible incluso cuando $k$
solo tiene dos elementos, pues entonces estos dos valores considerados
de $j_o$ coinciden. Así, la imagen $j_o$ de $j$ en $k(s')$ es distinta
de $0$ y de $12^3$, y se sabe que existe entonces una curva elíptica
sobre $S'$ con invariante $j$ igual a $j$.

\medskip
\noindent
6.3. Volvamos al caso en que $S$ solo se supone henseliano. Por analogía
con las definiciones de 4.1 para $\ell\neq p$, se puede definir entonces
una filtración
\begin{equation}\tag{6.3.1}
 T_p(A_K)\supset T_p(A_K)^{\mathrm{ef}}
 \supset T_p(A_K)^{\mathrm{et}}\supset0
\end{equation}
por subprogrupos estrictos de Barsotti--Tate, llamados respectivamente
la \underline{parte esencialmente fija} y la
\underline{parte esencialmente tórica} de $T_p(A_K)$.

Para ello, elijamos una extensión de Galois finita $K'$ de $K$ tal que
$A_{K'}$ tenga reducción semiestable sobre la normalización $S'$ de $S$
en $K'$. Queda entonces definida la parte fija $T_p(A_{K'})^f$ de
$T_p(A_{K'})=(T_p(A_K))_{K'}$ (2.2.3), así como la parte tórica (6.1).
Esta última es además el complemento ortogonal de la parte fija
$T_p(A'_{K'})^f$ en $T_p(A'_{K'})$ (5.2). Evidentemente,
\[
 T_p(A_{K'})^t\subset T_p(A_{K'})^f
\]
(esto se deduce de 5.3, o directamente de la construcción de la parte
tórica en 6.1), y el grupo de Galois de $K'/K$ preserva estos dos
subprogrupos estrictos de $(T_p(A_K))_{K'}$.

%%%%%%%%%%  scan idx 408  |  folio 397  %%%%%%%%%%

El descenso da, pues, la filtración buscada (6.3.1). Esta es
independiente de la elección de $K'$, por la propiedad de reducción
semiestable 3.3. El descenso da asimismo de manera evidente la propiedad
de ortogonalidad
\begin{equation}\tag{6.3.2}
 \begin{split}
 T_p(A_K)^{\mathrm{et}}
   &=\bigl(T_p(\sisi{A'_k}{A'_K})^{\mathrm{ef}}\bigr)^\perp,\\
 T_p(A_K)^{\mathrm{ef}}
   &=\bigl(T_p(\sisi{A'_k}{A'_K})^{\mathrm{et}}\bigr)^\perp.
 \end{split}
\end{equation}

Se podría dar una caracterización intrínseca de la filtración (6.3.1)
en términos del progrupo de Barsotti--Tate $X=T_p(A_K)$ sobre el cuerpo
de fracciones $K$ de $S$, en el espíritu de 5.8---al menos, por ahora,
en características desiguales. La parte esencialmente fija es el mayor
subprogrupo de Barsotti--Tate con ``buena reducción virtual sobre $S$'':
tras una extensión separable finita $K'$ de $K$, da un progrupo de
Barsotti--Tate $X_{K'}$ con buena reducción sobre la normalización $S'$
de $S$ en $\sisi{K}{K'}$, es decir, que se prolonga a un progrupo de
Barsotti--Tate sobre $S'$. Tal prolongación es necesariamente única
salvo isomorfismo único, por el teorema de Tate ya citado [33, Th. 4].
La parte esencialmente tórica se define entonces a partir de la parte
esencialmente fija por ortogonalidad (6.3.2).

Consignemos también la siguiente consecuencia inmediata de 5.8 y de
las definiciones; para $\ell=p$ se reduce a
\sisi{6.5 a)}{5.13 b)}.

\medskip
\noindent
\underline{Proposición} 6.3.3. \underline{Sea $\ell$ un número primo.
Si $\ell=p$, supongamos que $\operatorname{car}K=0$.}
\footnote{Esta hipótesis probablemente sea innecesaria; cf. 5.11.}
\underline{Entonces $A_K$ tiene buena reducción virtual sobre $S$
\sisi{(5.12 a))}{(5.13 b))} si y solo si su parte esencialmente fija}
$T_\ell(A_K)^{\mathrm{ef}}$
\underline{(4.1.1 y 6.3.1) es todo $T_\ell(A_K)$.}

%%%%%%%%%%  scan idx 409  |  folio 398  %%%%%%%%%%

\medskip
\noindent
6.4. Recordemos que, cuando $\ell\neq p$, en 4.3 a) se puso de
manifiesto una notable función de valores enteros sobre el grupo de
inercia $I$ de $S$, independiente de la elección de $\ell$. Cuando el
cuerpo residual es perfecto, el conocimiento de esta función basta
para reconstruir el exponente del conductor de $A$ (4.5): se toma,
sobre un cociente suficientemente grande $\Gamma=I/I'$ de $I$, el
producto interno de esta función con el conductor de Artin sobre
$\Gamma$ (cf. (4.6.2)). Conviene observar que esta función de valores
enteros sobre $I$ también puede recuperarse a partir del conocimiento
de $T_\ell(A_K)$ para $\ell=p$---al menos, por ahora, en características
desiguales.

Cuando el cuerpo residual $k$ es finito, una observación análoga se
aplica a la función de valores en $\mathbb Q$ considerada en 4.3 b),
definida sobre el subgrupo $\pi'$ de $\pi$ formado por los elementos
cuya imagen en $\pi_o$ tiene la forma $f_q^{\,n}$,
$n\in\mathbb Z$. Indicamos el principio de la verificación. Usando la
filtración (6.3.1), nos reducimos a asociar una función adecuada (sobre
$I$, respectivamente sobre $\pi'$) a un grupo pro-$p$ de
Barsotti--Tate $Y$ sobre $K$ que tenga ``buena reducción virtual sobre
$S$'' en el sentido de 5.12. Aquí $Y$ será uno de los tres cocientes
provenientes de la filtración (6.3.1).

En el caso de 4.3 a), una extensión conveniente del cuerpo residual en
$S$ (EGA $0_{\mathrm{III}}$ 10.3.1) nos reduce al caso de un cuerpo
residual perfecto, o incluso algebraicamente cerrado si se desea.
Elijamos una extensión de Galois finita $K'$ de $K$ tal que $Y_{K'}$
tenga buena reducción sobre $S'$. Tenemos entonces un grupo de
Barsotti--Tate $Y_{S'}$ sobre $S'$ y una acción de $G$ sobre $Y_{S'}$
compatible con su acción sobre $S'$. Si $k'$ es el cuerpo residual de
$S'$, entonces $G$ actúa, no necesariamente fielmente, sobre $k'$ (cuyo
cuerpo de invariantes es $k$), y sobre $Y_{k'}$ de manera compatible
con su acción sobre $k'$.

%%%%%%%%%%  scan idx 410  |  folio 399  %%%%%%%%%%

Por la teoría de Dieudonné [16 bis], el grupo pro-$p$ de
Barsotti--Tate $Y_{k'}$ tiene asociado canónicamente un módulo libre
finito $M$ (de rango igual al rango de $Y_{k'}$) sobre el anillo
$W(k')$ de vectores de Witt de $k'$. El grupo $G$ actúa evidentemente
sobre $M$ de manera compatible con sus acciones sobre $W(k')$, por
medio de sus acciones sobre $k'$. Por consiguiente, el subgrupo de
inercia $\Gamma$ de $G$ actúa $W(k')$-linealmente sobre $M$; como
$\Gamma$ es un cociente del grupo de inercia $I$, esto da una función
de traza sobre $\Gamma$, y por tanto sobre $I$, que es la función
buscada.

A priori, esta función toma valores en $W(k')$,
\footnote{De hecho, es inmediato que esta función toma siempre valores
en $W(k)$, pues es invariante bajo $\operatorname{Gal}(k'/k)$.}
pero coincide con la función análoga definida en términos de los
$T_\ell$ para $\ell\neq p$ siempre que la acción de $G$ sobre $Y_{S'}$
se obtenga de una acción de $G$ sobre un esquema en grupos
$p$-divisible liso $C'$ sobre $S'$, tomando $Y=T_p(C')$ (cf. [25 bis,
Th. 7, p. 737]). En este caso, por tanto, es de valores enteros. El caso
de 4.3 b) se trata de modo análogo: para definir una función de traza
sobre $\pi'$, se utiliza el método de la demostración de 4.3.2 para
``linealizar'' sobre $M$ la acción definida por un elemento $g\in\pi$
que está por encima de un $f_q^n$, con $n\geq0$.
